\subsection{Section summary}\label{subsec:doc-extraction-section-summary}

This section described the document-extraction pipeline, which turns PMC Open Access source papers into structured, provenance-preserving document elements. A layered design is used, because no single layout parser can recover scientific PDFs alone. Docling provides the structural backbone of our parser: a typed list of layout elements, and a section hierarchy read from heading-style blocks. The hybrid layout parser combines Docling's output with Microsoft's image-based Table Transformer (TATR) to recover tables pasted in as images. Then the detected figures and tables are cropped into separate images, and matched to captions with a geometry-based nearest caption heuristic. 

A two-pass extraction then masks invisible text, page headers, footers and sidebars, and detected figure and table crops, and re-extracts a layout that only contains body text and headings. This layout is assembled into hierarchical paragraphs: paragraphs split into multiple fragments across page- or column-breaks are merged, citation markers, URLs and unwanted text is removed, and each paragraph is tagged with the section path in which it appears. In the end, the final output is written into a relational database with stable identifiers so as to preserve provenance in later stages, and a set of on-disk file artifacts. 

The pipeline is designed with two commitments in mind. The first commitment is to preserve provenance so that each final rule could be traced back to its source. The second commitment is a bounded view of text extraction quality. The document-extraction pipeline is designed to produce elements that are useful for knowledge extraction downstream; it is not designed to reconstruct the PDF perfectly. The remaining failure modes are described in Subsection~\ref{subsec:doc-extraction-known-limitations}, and are quantified, where the rubric applies, in Chapter~\ref{chapter:Results}. 

The interface to the knowledge extraction pipeline, described in Subsection~\ref{subsec:doc-interface-to-knowledge-pipeline}, is narrower than the full stored representation of the documents. Although sections, section paths, figures, tables, and cross-references to these figures and tables are preserved and stored in the database, the downstream MAP stage only receives a stream of sentence-segmented text, with each sentence tagged with stable identifiers of the text elements it was derived from. Section~\ref{section:Knowledge_Extraction} builds on this provenance-preserving, sentence-level interface.